\section{The proof space as a graph}

\lean{RequestProject/ProofGraph.lean} reifies the spine of the development as a
\lean{SimpleGraph} on 25 nodes and proves its structural properties by \lean{decide}.
This section draws that graph.  Node colour is the tag, horizontal position is the rank
(= height in the DAG, \lean{Golf.Graph.rank}), and an arrow $b \to a$ means
``$a$ uses $b$'' (\lean{Golf.Graph.Uses}).

\begin{center}\footnotesize
\begin{tabular}{@{}llll@{}}
\toprule
\textcolor{tagatomic}{$\blacksquare$}~atomic & Mathlib, or one tactic & 10 nodes &
  \lean{card\_atomic}\\
\textcolor{tagreducible}{$\blacksquare$}~reducible & an iso/dual/rewrite collapses it &
  3 nodes & \\
\textcolor{tagglue}{$\blacksquare$}~glue & a small composition & 7 nodes & \\
\textcolor{tagstructural}{$\blacksquare$}~structural & introduces a new invariant &
  5 nodes & \lean{card\_structural}\\
\bottomrule
\end{tabular}
\end{center}

\begin{sidewaysfigure}
\centering
\resizebox{.96\textheight}{!}{%
\begin{tikzpicture}[x=1cm,y=1cm]
  % ---------------- rank 0 : atoms ----------------
  \node[atomic] (iSupMono)        at (0,10)   {iSup\_mono'};
  \node[atomic] (leAntisymm)      at (0,8.5)  {le\_antisymm};
  \node[atomic] (bddRow)          at (0,7)    {bddAbove\_range\_row};
  \node[atomic] (bddRowSups)      at (0,6)    {bddAbove\_range\_ciSup};
  \node[atomic] (cofinalRange)    at (0,4.5)  {isCofinalFor\_range\_iff};
  \node[atomic] (cofinalBdd)      at (0,3.5)  {IsCofinalFor.bddAbove};
  \node[atomic] (bddAdd)          at (0,2)    {bddAbove\_range\_add};
  \node[atomic] (csSupCofinal)    at (0,0.5)  {csSup\_eq\_csSup\_of\_\\forall\_exists\_le};
  \node[atomic] (colimitEqISup)   at (0,-1.2) {colimit\_eq\_iSup};
  \node[atomic] (finalDiag)       at (0,-2.4) {final\_diag\_of\_isFiltered};
  % ---------------- rank 1 ----------------
  \node[structural] (iSupPProd)   at (5,10)   {iSup\_pprod\\{\itshape currying}};
  \node[structural] (iSupCofinal) at (5,8.4)  {iSup\_eq\_iSup\_of\_cofinal\\{\itshape cofinality}};
  \node[structural] (ciSupPProd)  at (5,6.5)  {ciSup\_pprod\\{\itshape currying, bounded}};
  \node[glue] (ciSupCofinal)      at (5,5)    {ciSup\_eq\_ciSup\_of\_cofinal};
  \node[reducible] (bddUncurry)   at (5,3.9)  {bddAbove\_range\_uncurry};
  \node[structural] (ciSupLin)    at (5,0.5)  {ciSup\_eq\_ciSup\_of\_\\cofinal\_of\_linear};
  \node[structural] (finality)    at (5,-1.8) {iSup\_prod\_eq\_iSup\_\\diagonal\_functor};
  % ---------------- rank 2 ----------------
  \node[glue] (iSup2Cofinal)      at (10,9.2) {iSup₂\_eq\_iSup\_of\_cofinal};
  \node[glue] (ciSup2Diag)        at (10,5)   {ciSup₂\_eq\_ciSup\_diagonal};
  \node[glue] (bimonotone)        at (10,-1.8){iSup₂\_eq\_iSup\_diagonal\_\\of\_bimonotone};
  % ---------------- rank 3 ----------------
  \node[reducible] (iSup2Diag)    at (15,9.2) {iSup₂\_eq\_iSup\_diagonal};
  % ---------------- rank 4 ----------------
  \node[glue] (iSupOpDiag)        at (20,10.3){iSup\_op\_iSup\_diagonal};
  \node[reducible] (iInf2Diag)    at (20,8.1) {iInf₂\_eq\_iInf\_diagonal};
  % ---------------- rank 5 ----------------
  \node[glue] (ennrealAdd)        at (25,10.3){ennreal\_iSup\_add\_iSup};
  \node[glue] (cardinalAdd)       at (25,3.5) {cardinal\_ciSup\_add\_\\ciSup\_diagonal};
  % ---------------- rank guides ----------------
  \foreach \x/\r in {0/0,5/1,10/2,15/3,20/4,25/5}{
    \node[font=\small\itshape,text=black!45] at (\x,-3.6) {rank \r};}
  % ---------------- use edges (24) ----------------
  \draw[use] (leAntisymm) -- (iSupPProd);
  \draw[use] (leAntisymm) -- (iSupCofinal);
  \draw[use] (iSupMono)   -- (iSupCofinal);
  \draw[use] (leAntisymm) -- (ciSupPProd);
  \draw[use] (bddRow)     -- (ciSupPProd);
  \draw[use] (bddRowSups) -- (ciSupPProd);
  \draw[use] (leAntisymm) -- (ciSupCofinal);
  \draw[use] (csSupCofinal) -- (ciSupLin);
  \draw[use] (colimitEqISup) -- (finality);
  \draw[use] (finalDiag)  -- (finality);
  \draw[use] (cofinalRange) -- (bddUncurry);
  \draw[use] (cofinalBdd) -- (bddUncurry);
  \draw[use] (iSupPProd)  -- (iSup2Cofinal);
  \draw[use] (iSupCofinal)-- (iSup2Cofinal);
  \draw[use] (ciSupPProd) -- (ciSup2Diag);
  \draw[use] (ciSupCofinal) -- (ciSup2Diag);
  \draw[use] (bddUncurry) -- (ciSup2Diag);
  \draw[use] (finality)   -- (bimonotone);
  \draw[use] (iSup2Cofinal) -- (iSup2Diag);
  \draw[use] (iSup2Diag)  -- (iSupOpDiag);
  \draw[use] (iSup2Diag)  -- (iInf2Diag);
  \draw[use] (iSupOpDiag) -- (ennrealAdd);
  \draw[use] (ciSup2Diag) -- (cardinalAdd);
  \draw[use,rounded corners=6pt] (bddAdd) -- (2.2,2) -- (22.6,2) -- (cardinalAdd);
  % ---------------- bridge edges (2) ----------------
  \draw[bridge] (bimonotone) to[out=0,in=-80,looseness=.85]
    node[font=\small,text=bridgecol,pos=.45,sloped,above]{bridge} (iSup2Diag);
  \draw[bridge,rounded corners=8pt] (ciSupLin) -- (2.55,0.5) --
    node[font=\small,text=bridgecol,pos=.5,sloped,above]{bridge} (2.55,5) --
    (ciSupCofinal);
\end{tikzpicture}}
\caption{The dependency DAG of the development (\lean{Golf.Graph.G}), with the two
bridge links of \lean{Golf.Graph.G'} dashed in purple.  \lean{dep\_rank\_lt} states that
every solid arrow goes strictly left-to-right; \lean{atomic\_iff\_leaf} that the grey
nodes are exactly those with no incoming arrow; \lean{deps\_length\_le\_three} that no
node has more than three incoming arrows; \lean{card\_edges} that there are 24 solid
edges and \lean{card\_bridged\_edges} that the bridges cost exactly two more.}
\end{sidewaysfigure}

\subsection{Connectivity: a three-colouring, and the two bridges that kill it}

The interesting finding of \lean{ProofGraph.lean} is negative: the use-graph is
\emph{disconnected}.  The categorical route and the linear-order route reach the same
theorems through disjoint material.  This is proved by a graph colouring invariant
(\lean{colour\_eq\_of\_adj}, checked edge by edge with \lean{decide}, then propagated
along walks in \lean{colour\_eq\_of\_reachable}).

\begin{figure}[ht]
\centering
\begin{tikzpicture}[x=1cm,y=1cm,font=\footnotesize]
  \begin{scope}
    \draw[rounded corners=10pt,fill=tagglue!8,draw=tagglue!60,line width=.8pt]
      (0,0) rectangle (5.6,4.3);
    \node[font=\small\bfseries,text=tagglue] at (2.8,3.85) {colour 0 -- 19 nodes};
    \node[align=center] at (2.8,2.3)
      {the order-theoretic core:\\ currying, cofinality, boundedness,\\
       and \emph{both} headline theorems};
    \node[font=\scriptsize,text=black!60,align=center] at (2.8,0.6)
      {contains the root \lean{le\_antisymm}};
  \end{scope}
  \begin{scope}[xshift=6.3cm]
    \draw[rounded corners=10pt,fill=tagreducible!8,draw=tagreducible!60,line width=.8pt]
      (0,0) rectangle (4.4,4.3);
    \node[font=\small\bfseries,text=tagreducible] at (2.2,3.85) {colour 1 -- 2 nodes};
    \node[align=center] at (2.2,2.3)
      {\lean{csSupCofinal}\\ $\downarrow$\\ \lean{ciSupCofinalLinear}};
    \node[font=\scriptsize,text=black!60] at (2.2,0.6) {the linear-order route};
  \end{scope}
  \begin{scope}[xshift=11.4cm]
    \draw[rounded corners=10pt,fill=tagstructural!8,draw=tagstructural!60,line width=.8pt]
      (0,0) rectangle (5.0,4.3);
    \node[font=\small\bfseries,text=tagstructural] at (2.5,3.85) {colour 2 -- 4 nodes};
    \node[align=center] at (2.5,2.1)
      {\lean{colimitEqISup}, \lean{finalDiag}\\ $\downarrow$\\
       \lean{finalityCollapse}\\ $\downarrow$\\ \lean{bimonotoneDiag}};
    \node[font=\scriptsize,text=black!60] at (2.5,0.6) {the categorical route};
  \end{scope}
  \draw[bridge,bend left=20] (8.5,4.4) to
    node[above,font=\scriptsize,text=bridgecol]{bridge 2} (4.3,4.4);
  \draw[bridge,bend right=13] (13.9,-0.15) to
    node[below,font=\scriptsize,text=bridgecol]{bridge 1} (4.3,-0.15);
\end{tikzpicture}
\caption{\lean{not\_connected}: the colouring is constant on connected components and
takes three values, so \lean{G} is disconnected --- the categorical proof of the collapse
shares \emph{no} lemma with the order-theoretic one.  Adding the two purple bridge
edges makes the space connected (\lean{bridged\_connected}).}
\end{figure}

\subsection{Why connectivity is proved by descent, not by search}

Reachability is expensive for the kernel, so the positive statement is reduced to a
local one: a parent function \lean{toRoot} together with a strictly decreasing
\lean{depth}.  This is a spanning tree of the bridged graph, rooted at
\lean{le\_antisymm}.

\begin{sidewaysfigure}
\centering
\resizebox{.96\textheight}{!}{%
\begin{tikzpicture}[x=1cm,y=1cm,font=\scriptsize]
  \tikzset{tn/.style={rounded corners=2pt,draw=black!45,fill=white,
    font=\ttfamily\scriptsize,inner sep=2.5pt,align=center}}
  \node[tn,fill=tagatomic!18,draw=black!65] (r) at (0,0) {le\_antisymm};
  % depth 1
  \node[tn] (a1) at (5.2, 5.5) {iSup\_pprod};
  \node[tn] (a2) at (5.2, 2.5) {iSup\_eq\_iSup\_of\_cofinal};
  \node[tn] (a3) at (5.2,-1.0) {ciSup\_pprod};
  \node[tn] (a4) at (5.2,-4.0) {ciSup\_eq\_ciSup\_of\_cofinal};
  % depth 2
  \node[tn] (b1) at (10.9, 4.0) {iSup\_mono'};
  \node[tn] (b2) at (10.9, 1.6) {iSup₂\_eq\_iSup\_of\_cofinal};
  \node[tn] (b3) at (10.9,-0.4) {bddAbove\_range\_row};
  \node[tn] (b4) at (10.9,-1.8) {bddAbove\_range\_ciSup};
  \node[tn] (b5) at (10.9,-3.6) {ciSup\_eq\_ciSup\_of\_\\cofinal\_of\_linear};
  \node[tn] (b6) at (10.9,-5.8) {ciSup₂\_eq\_ciSup\_diagonal};
  % depth 3
  \node[tn] (c1) at (16.9, 1.6) {iSup₂\_eq\_iSup\_diagonal};
  \node[tn] (c2) at (16.9,-3.6) {csSup\_eq\_csSup\_of\_\\forall\_exists\_le};
  \node[tn] (c3) at (16.9,-5.6) {bddAbove\_range\_uncurry};
  \node[tn] (c4) at (16.9,-7.4) {cardinal\_ciSup\_add\_\\ciSup\_diagonal};
  % depth 4
  \node[tn] (d1) at (23.0, 3.4) {iInf₂\_eq\_iInf\_diagonal};
  \node[tn] (d2) at (23.0, 1.6) {iSup\_op\_iSup\_diagonal};
  \node[tn] (d3) at (23.0,-0.2) {iSup₂\_eq\_iSup\_diagonal\_\\of\_bimonotone};
  \node[tn] (d4) at (23.0,-5.0) {isCofinalFor\_range\_iff};
  \node[tn] (d5) at (23.0,-6.4) {IsCofinalFor.bddAbove};
  \node[tn] (d6) at (23.0,-8.0) {bddAbove\_range\_add};
  % depth 5
  \node[tn] (e1) at (29.2, 1.6) {ennreal\_iSup\_add\_iSup};
  \node[tn] (e2) at (29.2,-0.2) {iSup\_prod\_eq\_iSup\_\\diagonal\_functor};
  % depth 6
  \node[tn] (f1) at (35.0, 0.6) {colimit\_eq\_iSup};
  \node[tn] (f2) at (35.0,-1.0) {final\_diag\_of\_isFiltered};
  \foreach \s/\t in {a1/r,a2/r,a3/r,a4/r,b1/a2,b2/a2,b3/a3,b4/a3,b5/a4,b6/a4,
                     c1/b2,c2/b5,c3/b6,c4/b6,d1/c1,d2/c1,d3/c1,d4/c3,d5/c3,d6/c4,
                     e1/d2,e2/d3,f1/e2,f2/e2}
    {\draw[-{Stealth[length=2mm]},draw=black!50] (\s) -- (\t);}
  \foreach \d/\x in {0/0,1/5.2,2/10.9,3/16.9,4/23.0,5/29.2,6/35.0}{
    \node[text=black!45,font=\scriptsize\itshape] at (\x,-9.4) {depth \d};}
\end{tikzpicture}}
\caption{The spanning tree given by \lean{toRoot}: every node other than the root is
adjacent in \lean{G'} to its parent (\lean{adj\_toRoot}) and strictly closer to the root
(\lean{depth\_toRoot\_lt}).  Strong induction on \lean{depth} then gives
\lean{reachable\_root} and hence \lean{bridged\_connected}.  Note that the tree uses
\emph{both} bridge edges, which is why it does not exist in \lean{G}.}
\end{sidewaysfigure}

\subsection{Cofinality as a bipartite graph}

The hypothesis pair $h_1 : \forall i, \exists j,\ f_i \le g_j$ and
$h_2 : \forall j, \exists i,\ g_j \le f_i$ of
\lean{Golf.iSup\_eq\_iSup\_of\_cofinal} is a combinatorial condition on a bipartite
graph: every vertex on each side has at least one out-neighbour on the other.

\begin{figure}[ht]
\centering
\begin{tikzpicture}[x=1.9cm,y=1.55cm,font=\footnotesize]
  \foreach \i in {1,...,5}{\node[latticept,fill=tagstructural!30] (f\i) at (\i,1.4) {};
    \node[font=\scriptsize,above=2pt] at (f\i) {$f_\i$};}
  \foreach \j in {1,...,5}{\node[latticept,fill=tagreducible!30] (g\j) at (\j,-0.8) {};
    \node[font=\scriptsize,below=2pt] at (g\j) {$g_\j$};}
  \foreach \i/\j in {1/1,2/2,3/2,4/4,5/5}{
    \draw[-{Stealth[length=1.8mm]},tagstructural] (f\i) -- (g\j);}
  \foreach \j/\i in {1/2,2/3,3/3,4/5,5/5}{
    \draw[-{Stealth[length=1.8mm]},tagreducible,dashed] (g\j) -- (f\i);}
  \node[anchor=west,align=left] at (6.0,1.4) {$h_1$: solid arrows,\\ $f_i \le g_{j}$};
  \node[anchor=west,align=left] at (6.0,-0.8) {$h_2$: dashed arrows,\\ $g_j \le f_{i}$};
  \node[anchor=west,align=left,text=bridgecol] at (6.0,0.3)
    {no perfect matching is required:\\ only that no vertex is isolated};
\end{tikzpicture}
\caption{Mutual cofinality.  The proof never chooses a canonical partner: it applies
\lean{iSup\_mono'} twice, which consumes exactly the existence of \emph{some}
out-neighbour.}
\end{figure}

\clearpage
